\documentclass[journal]{IEEEtran}
\usepackage{graphicx}
\usepackage{amsmath}
\usepackage{booktabs}
\usepackage{url}
\usepackage[utf8]{{inputenc}}


\title{Lane Perception for Automated Driving with Human-Anchored, Engine-Certified Ground Truth: An Explicit Counting Contract and Pre-Registered Single-Factor Evaluation}
\author{Zheyu Yuan (袁哲宇)}

\begin{document}
\maketitle
\begin{abstract}
Lane perception for automated driving is bottlenecked by two conflated problems: the cost of pixel-accurate annotation, and the absence of an agreed definition of a "lane line". We present a closed-loop research programme that attacks both. Every label carries an explicit provenance credential --- human revision, agent review, or engine-certified truth --- so no claim has to guess what truth it rests on, and the reported rounds added no manual annotation. Evaluation is made explicit and versioned: a \emph{counting contract} gives every ratio an explicit numerator and denominator, and a \emph{protocol} fixes the geometric scope of a line candidate, with a hash binding the sealed set to its definition. One segmentation arm (six seeds) is evaluated on two reference sets, a road-disjoint held-out set and in closed-loop driving. Three findings follow. (i) With the weights fixed, changing the post-processing switches and the evaluation scope changes the readings and can reverse the ranking of two arms; about 30\% of annotated line pixels fail the appearance predicate the adopted definition uses. (ii) The one-shot confirmation does \textbf{not} reproduce the acceptance numbers (label-scope recall 0.460, precision 0.763) --- the headline of that independent test. (iii) The same definition change came with no observed centre-line crossing in four runs, a 66\% drop in lane-source availability, and an acceptance gate still failing 0/4; the blocker is a measured static body-sweep deadlock. All changes were pre-registered and judged by rules fixed before the runs.
\end{abstract}

\section*{Highlights}
- Ground-truth provenance is credentialed and enforced: human revision, agent review, engine-certified.
- A counting contract gives every reported ratio an explicit numerator and denominator.
- The protocol hash binds a sealed held-out set to the definition it is read under.
- The two natural definitions of a lane line fail different gates on identical data.
- A one-shot held-out confirmation does not reproduce the development-set numbers.
- Five pre-registered driving factors were rejected; the blocker is a measured deadlock.
\section{Introduction}
Lane-keeping is the lowest-level competence of an automated driving stack, and it is still
surprisingly hard to evaluate. Two obstacles recur in the literature. The first is annotation cost:
pixel-accurate lane labels are expensive, so public datasets are small relative to the diversity of
road appearance, and every new operational domain demands a new labelling campaign. The second is
definitional: papers report "lane detection" accuracy against ground truth that may include or exclude
faded, occluded, dashed, yellow, or shadowed markings, and may or may not count markings beyond the
ego lane. Two systems can therefore disagree by tens of accuracy points while disagreeing about
nothing except the definition of the object they detect.
This paper reports a research programme built around those two obstacles. It is not a new network
architecture; it is an attempt to make the \emph{data} and the \emph{evaluation} rigorous enough that
architectural progress can be measured. The programme has three components.
\textbf{Credentialed truth provenance.} Every label in this programme carries a provenance credential, and the credential is
enforced rather than declared: a command-line claim cannot promote a label (Section III-B). The ladder has three rungs --
human frame-by-frame revision, per-frame agent review, and engine-derived truth -- and the engine rung is admitted only
through four explicit certification gates: a credential
gate (the label source must be declared and verifiable), a paint-appearance gate (the pixels claimed
as line must look like paint in the image), a road-type gate (the frame must carry the road-type
column that makes pavement decidable), and a negative-evidence gate (frames claimed to contain no
line must show evidence of their own, not merely the absence of a label). The rounds reported here add no new manual annotation: the experimental variable is the composition of the
engine-certified increment on top of the human-revised anchor, and Section VII states what the anchor's small size costs
us.
\textbf{An explicit counting contract.} Every reported ratio is defined by a counting contract that names
its numerator and denominator over an explicit set chain: truth-confirmed paint frames, evaluated
perception candidates on those frames, candidates with a reference on their own side, matched
candidates, matches with decidable predicted and reference roles, and role-agreeing matches. Coverage, identity and
role are therefore instance-level ratios with stated denominators; precision and recall remain
pixel-level and are reported in two \emph{scopes} (label pixels, and paint-like pixels) because the two
answer different questions.
\textbf{Versioned protocol definitions with hash-gated evaluation.} A protocol fixes the geometric scope
of a line candidate and is a first-class, versioned artefact. A protocol change alters the evaluation
hash, and both the frozen held-out set and the one-shot confirmation record that hash, so a
confirmation produced under one definition cannot be silently attributed to another. This
mechanism caught a real defect during the work: before it was fixed, the hash did not follow the
protocol, and a set sealed under the old definition would have authorised a confirmation under the
new one.
\textbf{Method discipline.} Every change in this paper was pre-registered with its reading rule before
the runs, and each
pre-registration names the protocol version and the repository commit it was frozen against; the rules include safety non-inferiority read as a maximum over runs, an explicit refusal
to treat an unmeasured quantity as a pass, and a prohibition on relaxing the frozen gates. Changes
were applied one factor at a time and judged against alternating control/factor runs. Where a
pre-registered rule produced a verdict we considered too coarse, we report the verdict as given and
record the rule's defect rather than reinterpreting the result; one such defect --- safety metrics
that are frame \emph{counts} and therefore scale with how far the car travelled --- is analysed in
Section V-F and quantified in the supplement.
\textbf{Contributions.} (1) A credentialed truth-provenance ladder (human revision / agent review / engine-certified) with enforcement, plus a
certification-gated engine pipeline for the machine-generated increment; the composition of every pool is stated
explicitly, and the reported rounds add no new manual annotation. (2) The counting contract, and the demonstration that the two
natural readings of "line" fail different gates on the same data --- a definitional disagreement that moves
about 30\% of annotated line pixels out of the adopted scope (which retains 69\% of annotated line pixels on
the development pool and 89\% on the limited-class pool). (3) A protocol/version/hash mechanism that binds a held-out set to the
definition under which it is read. (4) A one-shot, road-disjoint confirmation with an explicit
UNKNOWN policy, whose negative result (no reproduction of the development-set numbers) we report as
the finding. (5) A measured localisation of why a perception stack that passes its offline gates
still fails a closed-loop acceptance gate, including a static-deadlock mechanism and the observation
that the "on-pavement" gate was in fact being driven by a neural obstacle layer rather than by
pavement evidence.
\section{Related Work}
\textbf{Lane detection methods.} Modern lane detection is dominated by deep segmentation and by
parametric or anchor-based curve regression. Encoder--decoder segmentation with dilated convolutions
and a discriminative loss (LaneNet \cite{neven2018lanenet}) and message-passing across rows (SCNN
\cite{pan2018scnn}) established the segmentation line of work that the present system follows;
cross-layer refinement with anchor-based heads (CLRNet \cite{zheng2022clrnet}) improved curve accuracy
and speed. Chinese-language work in the same family explores lightweight variants of CLRNet
\cite{chen2026clrnetlight}, B-spline end-to-end regression \cite{huang2026bspline}, graph/transformer hybrids
\cite{jia2026gnn}, transformer with cross-attention \cite{li2026xca}, parallel detection over image
sequences \cite{zhu2024parallel}, multi-scale feature aggregation \cite{hu2026multiscale}, DeepLabV3+
variants \cite{shuai2026deeplab} of the original architecture \cite{chen2018deeplab}, symmetry-prior and key-point methods for complex scenes
\cite{lu2026symmetry}, and 3-D lane estimation with convolutional architectures \cite{chen2025threed}.
These works share our object of study but differ in what they optimise: they report accuracy against
datasets whose annotation conventions differ from each other, which is exactly the ambiguity this
paper makes explicit.
\textbf{Datasets and ground truth.} TuSimple \cite{tusimple2017}, CULane \cite{pan2018scnn} and Cityscapes \cite{cordts2016cityscapes} remain the
standard benchmarks; their annotation conventions (lane presence per row, presence of far-field
markings, treatment of occlusion) differ, and neither reports the annotation's \emph{appearance}
properties. Ground-truth generation has been studied as a problem in its own right: Borkar et al.
\cite{borkar2012groundtruth} built an automated ground-truth process around time-slicing visualisation
for a night-time lane detector, and McCall and Trivedi \cite{mccall2006video} surveyed and evaluated
video-based lane estimation with an explicit eye on evaluation methodology. Our work continues that
line with two additions: truth derived from engine state rather than from imagery, and a counting
contract that makes the denominators of every reported ratio inspectable.
\textbf{Automatic labelling and simulation.} Self-training and auto-labelling are widely used to reduce
annotation cost, but their evaluation inherits the definitional problem above. Simulator-derived
labels are common in synthetic benchmarks; less common is the discipline of certifying each
generated label against image evidence, which we adopt here because a label that cannot be
recovered from the image is not a learnable target.
\textbf{Evaluation discipline.} Pre-registration is standard in empirical sciences and rare in perception
research. Our contribution here is not the concept but a worked, fully recorded instance: a
pre-registration, an alternating single-factor A/B, a fixed reading rule, and a verdict per factor ---
including five factors whose verdict was negative.
\section{System and Method}
\subsection{Overall architecture}
The stack is a perception-led driving system. A camera ring (eight mounts) and a LiDAR range image
feed a semantic segmentation head whose line output is converted to world-space markings; a lane
reference is selected from those markings under the protocol definition; a planner produces a path;
a safety monitor arbitrates the tick against hard checks (pose, boundaries, occupancy, obstacle
risk) and publishes a speed and level; the controllers consume the arbitrated path. The iron rule of
the stack is that the lateral reference comes from perception only --- the map answers "where to",
never "where the lane is". the supplement shows the candidate-scope pipeline with the measured revocation
rate discussed in Section V-E.
\subsection{Truth provenance, credentials, and certification gates}
Two dimensions must not be conflated here. \textbf{Provenance} says who or what produced a label (human
revision, agent review, engine); \textbf{certification} says whether that label passed the four gates below.
A human-revised frame can fail the appearance gate and an engine frame can pass all four, so admission is
decided by the gate result, never by the rung. With that separation stated, the programme distinguishes
three rungs of truth and records the rung with every label: \textbf{human revision} (a human
annotates or corrects the line labels frame by frame), \textbf{agent revision} (a per-frame agent review of engine output), and
\textbf{engine-certified} truth (the engine's own line geometry projected into the camera). Two pools are human-revised: the
training anchor (5 directories, 75 frames) and the development/evaluation packs (9 packages, 117 frames), which is why the
acceptance results of Section V-A are measured against human annotation. The machine-generated increment that the ablations
vary is engine-certified; it is admitted only when: (i) the credential is present and parseable, and the label source is
declared (frames whose provenance cannot be named are excluded rather than guessed); (ii) the pixels
claimed as line satisfy a paint-appearance predicate --- white paint requires high value, low chroma
and a red/blue balance; yellow paint requires a hue ratio band; the predicate is a single-source
module so that a threshold change is a versioned change, not a local edit; (iii) the frame carries the
road-type column, without which "on pavement" is undecidable; (iv) a frame claimed to contain no line
carries positive evidence of its own (a certified negative), so that "no label" is never read as
"no line". Frames failing any gate are excluded from both training and evaluation and are counted in
the audit, not silently dropped.
The gate bookkeeping is itself a result: the supplement shows the negative-pool funnel for one collection round (admitted, suspect, contains paint, undecided) --- a pool that looked plentiful on disk shrank to a certified subset before it was ever used as training signal.
\subsection{The counting contract}
The contract is stated over integer counters that are accumulated per frame; the ratios are computed once
from those integers, never by averaging per-directory ratios. $P$ counts frames whose annotation carries a
line on pavement (decided frame by frame: a frame is not admitted because another frame in its group has a
line); $C$ counts the perception candidates inside $P$ frames that enter the evaluation;
$C_{\mathrm{out}}$ counts candidates on frames outside $P$ and is reported separately; $R \subseteq C$
counts candidates whose own side has an available reference; $M \subseteq R$ those satisfying the frozen
matching condition; $L \subseteq M$ those where both the prediction and the reference have a decidable
left/right role; and $A \subseteq L$ those whose role agrees. The reported instance-level ratios are
coverage $|R|/|C|$, identity $|M|/|R|$ and role agreement $|A|/|L|$; pixel-level precision and recall are
reported in two scopes (label, paint). 1 shows the contract.
\begin{figure*}[!t]
\centering
\includegraphics[width=\textwidth]{../../logs/paper_build/figures/fig30_counting_contract.png}
\caption{Counting contract (schematic): every ratio has an explicit numerator and denominator, over the frozen counter definitions; coverage is candidate-reference availability, not detection recall.}
\label{fig:contract}
\end{figure*}
Two consequences have to be read off this definition rather than assumed. First, coverage is the
\textbf{candidate-reference availability}: the share of evaluated candidates whose own side has truth to compare
against, not the share of true markings that were detected. A line-bearing frame on which the perception
produces no candidate contributes nothing to $|R|/|C|$; its miss is carried by pixel-level recall instead.
Second, candidates on frames outside $P$ are counted as $C_{\mathrm{out}}$ and reported, so that "no
reference available" is never scored as a coverage failure. We therefore claim no instance-level recall:
the three conditional ratios above cannot substitute for one, and reporting one would require the total
number of ground-truth instances together with an explicit one-to-one assignment rule. Every ratio still
has an explicit denominator, and the frame-level counter $P$ is reported separately from the candidate-level
counters, so a run with few candidate-bearing frames cannot masquerade as a good one.
\subsection{Protocol definitions}
A protocol fixes three things about a line candidate: its \textbf{lateral scope} (candidates farther than
a threshold from the ego's own lane are not "my lane's lines"), whether **same-side near-duplicate
markings are merged\textbf{ into one instance, and whether an }appearance gate** removes candidates whose
pixels do not look like paint. The frozen baseline protocol applies none of the three (lateral scope
off, merging off, appearance off); the adopted protocol applies all three with a lateral threshold of
5.5 m, a merge radius of 1.0 m and the appearance predicate of Section III-B. the supplement shows the
switch matrix and the live hash; the hash is a function of the active protocol, so evaluation
artefacts cannot be reattributed across definitions.
The lateral scope is not arbitrary: it is set from the measured geometry of the reference pool, shown in the supplement (261 042 marking pixels, median offset −0.17 m, median absolute offset 1.22 m). A scope of 5.5 m retains the own lane's markings and their immediate neighbours while excluding the far-field markings that dominate the false pairings discussed in Section V-E.
\subsection{Training-composition design space}
Because ground truth is generated rather than annotated, the composition of the training pool is a
first-class experimental variable. We study it along one axis: the \textbf{negative dose}, i.e. how many
certified line-free frames are mixed into the pool relative to the base composition, together with
composition variants that place markings in specific geometric relations (near/far pairs, same-side
pairs). Section V-C reports the resulting dose--response.
\subsection{Held-out final set: sealing, hash gating, one-shot consumption}
An independent confirmation set is composed under three constraints: it must be \textbf{road-disjoint}
from every training pool (identity taken from the generation metadata, not from directory names), it
must be certified by the gates of Section III-B, and its content must be frozen by a digest over
per-frame hashes. Sealing records the frame digests, the protocol hash and the composition; access
is then gated on the protocol hash, and a single confirmation \textbf{consumes} the set: after a
successful confirmation the set may not be read again, and a refused access is recorded in a ledger
without consuming it. The confirmation program evaluates the candidate on the sealed frames and
records which quantities were measured; quantities that were not measured are recorded as UNKNOWN
and may not be reported as confirmed. The full flow, with the live seal, is shown in the supplement.
\subsection{Driving stack and safety arbitration}
In closed loop the same perception feeds a planner and a safety monitor whose arbitration ladder
evaluates rules in a fixed order; soft rules accumulate speed caps, while a rule whose worst level is
`minimal_risk` stops the tick. the supplement lists the ladder. Two rules matter for Section V-E: the
\textbf{planned body-sweep rule}, which stops the tick when the body sweep of the planned path would cross
a lane boundary closer than a hard threshold (4 m), and the \textbf{on-pavement gate}, which revokes the
perception lane when the lane centre does not lie on observed drivable surface.
\section{Experimental Setup}
\textbf{Reference sets.} Two disjoint sets are used for the primary acceptance: a development pool of 9
certified scenes (117 frames) and a limited-class pool of certified scenes chosen for near-field
background structure (24 frames per arm). Neither shares frames with training pools.
The limited-class pool is selected, not sampled: candidate scenes are ranked by a model-independent near-field structure-density criterion computed from the human annotation (the supplement), so the pool concentrates on the scenes where identity and role are hardest rather than on the average scene. A third,
road-disjoint set of 200 frames (96 line-bearing) is sealed for the one-shot confirmation of
Section V-D.
Disjointness is audited with pose, not with names: every candidate collection is checked against the development anchors under a 50 m buffer and gets a recorded verdict (supplement). The same machinery rejected the first two compositions of the sealed set, which is why the sealed composition also carries a content-level audit (the supplement).
\textbf{Arms.} All segmentation arms share one initialisation, one budget and one recipe; they differ only in training-pool
composition (the dose arms) or in a single post-processing factor (protocol components). Composition is explicit: a
human-revised anchor of 5 directories (75 frames) plus a varying number of engine-certified packages (2 paired-line
packages of 16 frames and 12 certified negative packages of 96 frames in the reported base), with the negative dose as the
varied axis. Six seeds are used for acceptance; composition arms are reported with their seed counts.
\textbf{Metrics and gates.} The frozen gates are coverage $\ge 0.80$, identity $\ge 0.60$, precision
$\ge 0.40$, recall $\ge 0.70$ and role agreement $\ge 0.70$; recall is read in the paint scope for
the adopted protocol and in the label scope for the baseline, and both readings are always reported.
The verdict rule changed between the two rounds and the change is part of the record rather than a detail:
the baseline round's stored verdicts require \textbf{every seed} to pass, while the adopted round stores
`verdict_mode = mean` and reports the per-seed values alongside. Section V-A therefore gives the stored
verdicts and a re-tabulation of the same per-seed numbers under each rule, so that no gate outcome has to
be read under an unstated rule.
\textbf{Driving experiments.} Closed-loop runs use one fixed scenario (a town link), alternating control
and factor arms, four runs per arm; hard targets are the benchmark's fixed checks
(frames, reversing, centre and edge crossings by reference line \emph{and} by body sweep, stall,
on-road, collisions, goal). Reading rules were fixed before each round: hard gates require 4/4,
safety items are compared as the maximum over runs, availability as median with range, and an
UNKNOWN never releases a gate.
\textbf{Discipline.} Every factor in Section V was pre-registered with its reading rule; no threshold was
relaxed at any point;
Timing hygiene is recorded rather than estimated: the supplement is the repeated-measurement protocol used before any latency claim (quiet-machine precondition, repeated runs, smaller value taken --- a leaked simulator instance once inflated p95 by a factor of 7.6). every verdict file, manifest and per-frame telemetry trace is retained, and
every figure in this paper is regenerated from those files by two scripts.
\section{Results}
\subsection{Acceptance under two definitions}
Table I reports the two definitions on both sets, and Table II re-tabulates the recall gate of the same
checkpoints under both post-processing settings and both scopes. Under the baseline definition the stored
verdict fails identity on \textbf{both} sets (0.433 dev, 0.414 limited-class) and additionally precision (0.402)
and role (0.674) on the limited-class set; under the adopted definition all five gates pass on both sets by
the mean rule (coverage 0.956/0.861, identity 0.665/0.716, precision 0.894/0.659, paint recall 0.812/0.885,
role 0.883/0.793). 2 shows the same numbers against the thresholds.
\begin{figure*}[!t]
\centering
\includegraphics[width=\textwidth]{../../logs/paper_build/figures/fig1_gate_matrix.png}
\caption{Gate metrics against the frozen thresholds for the two protocol definitions, development pool (a) and limited-class pool (b).}
\label{fig:gates}
\end{figure*}
Three qualifications belong with this table. (i) The two rows were not judged by the same rule. Re-tabulated
under a common mean rule, the baseline's limited-class precision (0.4018) is above its 0.40 threshold even
though 4 of 6 seeds fall below it, and the baseline's development recall (0.736) is above its 0.70 threshold
even though 2 of 6 seeds fall below it; under a common all-seeds rule the adopted definition's development
label-scope recall fails 6 of 6 seeds while its paint-scope recall fails 1 of 6. (ii) The protocol change
moves three post-processing switches and the evaluation scope at the same time, so this table cannot
attribute the whole difference to the definition; Table II separates the scope from the post-processing for
recall, and the per-switch single-factor results are in Section V-B. (iii) The candidate-level gates are the
stable part of the picture: the identity gate is missed by 6 of 6 baseline seeds on both sets and role by 5
of 6 on the limited-class set, while under the adopted definition no seed misses any of the three
(the supplement).
the supplement plots the limited-class pool seed by seed under both definitions, with each candidate-level
gate marked when missed; it is drawn from the same two verdict files as Table I, so the figure and the table
cannot disagree about which run produced them.
\begin{table*}[!t]
\centering
\caption{The two protocol definitions on both reference sets, as stored by each round under its own
verdict rule (means over six seeds; bold marks a gate the stored verdict failed; the adopted protocol is read
in paint scope, the baseline in label scope). The rule differs between the two rows --- see the text and
Table II.}
\begin{tabular}{@{}lllllll@{}}
\toprule
Set & Definition & Coverage & Identity & Precision & Recall & Role \\
\midrule
Dev (9 scenes, 117 frames) & baseline (label scope) & 0.952 & \textbf{0.433} & 0.610 & 0.736 & 0.868 \\
Dev & adopted (paint scope) & 0.956 & 0.665 & 0.894 & 0.812 & 0.883 \\
Limited-class (24 frames) & baseline (label scope) & 0.819 & \textbf{0.414} & \textbf{0.402} & 0.804 & \textbf{0.674} \\
Limited-class & adopted (paint scope) & 0.861 & 0.716 & 0.659 & 0.885 & 0.793 \\
\bottomrule
\end{tabular}
\end{table*}
\begin{table}[!t]
\centering
\caption{The recall gate re-tabulated: the same six checkpoints per pool, read under both
post-processing settings and both scopes, with the number of seeds below the 0.70 threshold. The paint scope
reads each checkpoint under the adopted definition; the label scope under the baseline one.}
\begin{tabular}{@{}lllll@{}}
\toprule
Pool & Post-processing & Scope & Mean recall & Seeds below 0.70 \\
\midrule
Development & baseline switches (v7) & label & 0.736 & 2/6 \\
Development & baseline switches (v7) & paint & 0.813 & 1/6 \\
Development & adopted switches (v8) & label & 0.559 & 6/6 \\
Development & adopted switches (v8) & paint & 0.812 & 1/6 \\
Limited-class & baseline switches (v7) & label & 0.804 & 0/6 \\
Limited-class & baseline switches (v7) & paint & 0.891 & 0/6 \\
Limited-class & adopted switches (v8) & label & 0.775 & 0/6 \\
Limited-class & adopted switches (v8) & paint & 0.885 & 0/6 \\
\bottomrule
\end{tabular}
\end{table}
\subsection{The disagreement is definitional}
Two things move together here and they have to be separated. The \textbf{scope} changes the reading: Table II
shows that under a fixed post-processing setting, moving from the label scope to the paint scope changes
development recall from 0.736 to 0.813 (baseline switches) and from 0.559 to 0.812 (adopted switches) --- in
the second case a 25-point swing with no change to the weights. The \textbf{post-processing} also changes the reading, and not in the same direction: with the scope held at label, the adopted switches cost 18 points of recall (0.736 $\rightarrow$ 0.559). Table V, which turns each switch on alone, attributes that cost entirely to the appearance gate: the lateral scope and the merge leave label recall untouched (0.736 in both) and act on identity instead (0.432 $\rightarrow$ 0.621 and 0.432 $\rightarrow$ 0.459 respectively), while the appearance gate is the only switch that trades label recall for pixel precision. What the scope change is worth in pixels is measured, not assumed:
the paint scope retains 69\% of annotated line pixels on the development pool and 89\% on the limited-class
pool, and the supplement shows the per-seed label-scope and paint-scope recall with the non-paint fraction of
label line pixels overlaid (\textasciitilde{}0.30). Those excluded pixels are the ones the appearance predicate does not
select; Section VI states what that does and does not license us to claim.
The instance-level gate is more sensitive than that, and in a way that matters for model selection.
the supplement reads 20 checkpoints (two arms, ten seeds each) under the label-scope and surface-scope
identity definitions: the base arm moves from a mean of 0.399 to 0.867 and the 6$\times$ arm from 0.429 to 0.771.
The absolute values nearly double for one arm and fall for the other, and **the ranking of the two arms
reverses** (base below 6$\times$ by 0.031 under the label scope, above it by 0.095 under the surface scope). The
same reversal is recorded in the repository's protocol proposal for the surface scope, which is why that
scope is reported rather than promoted: a gate that can change which checkpoint is selected is a
model-selection decision, not a reporting detail. An identity threshold must therefore name its scope,
exactly as the recall gate does. This experiment uses its own pool and its own ten seeds per arm, so its
numbers are not comparable with Table I's six-seed means and are not presented as such.
The effect is not specific to the delivered arm. On a second arm trained on a different composition (near/far pair packages, six seeds, same frozen protocol), moving from the label to the paint scope raises recall from 0.459 to 0.673 (+0.214) against +0.252 for the delivered arm, and the identity gate reads 0.658 against 0.665. The direction of the scope effect reproduces on the second model; its magnitude is arm-specific, which is why this paper reports the scope with every number instead of quoting a single calibration.
Two further single-factor results constrain how the adopted definition behaves. the supplement (boundary map)
sweeps mask dilation: dilation buys recall in both sets but costs precision, and the limited-class set pays
much more for it (at one pixel of dilation its precision falls from 0.65 to 0.39, a drop of 0.26, against
0.89 to 0.84, a drop of 0.05, on the development pool; the two rows are read at their own dilation
settings), so no single dilation value clears both sets. 3 shows the appearance
gate acting exactly as designed: precision rises sharply (dev 0.650 $\rightarrow$ 0.884; limited-class
0.387 $\rightarrow$ 0.650) and IoU improves on the limited-class set (0.347 $\rightarrow$ 0.532) at the cost of recall
(dev 0.766 $\rightarrow$ 0.595).
\begin{figure*}[!t]
\centering
\includegraphics[width=\textwidth]{../../logs/paper_build/figures/fig14_appearance_gate.png}
\caption{Appearance-gate ablation on the pixel metrics (line IoU, precision, recall) for both sets, seed 42 --- the only switch that trades label recall for precision (Table V).}
\label{fig:appearance}
\end{figure*}
The appearance predicate is a mask gate, not a truth criterion: its own module documents 13k--68k px of hits
on certified line-free frames (bright pavement and bright gravel both satisfy the white branch), which is why
it is used to remove candidates rather than to declare where lines are. With that caveat, the instance-level
probe records high annotation coverage for the observed limited-class reference instances, with
some values below 1.0. This finite sample does not establish complete truth for every surviving
candidate. Same-side parallel merging is a separate post-processing factor with a small identity
effect in Table V; its radius scan must not be confused with the lateral-scope scan. The two mask
thresholds were swept as well; on this pool the negative side is not measurable (no eligible
negative frames), so that sweep is reported in the supplement rather than plotted as a zero.
\begin{table*}[!t]
\centering
\caption{Post-processing single factors on the frozen development pool, one switch at a time (six seeds each; C0 is the baseline definition, C4 the adopted one). The lateral scope is an identity lever (identity 0.432 $\rightarrow$ 0.621, all six seeds positive), the appearance gate is a precision lever that pays label recall (0.610 $\rightarrow$ 0.894 precision, 0.736 $\rightarrow$ 0.559 recall) and does not raise identity, and the merge is a small identity effect (+0.027).}
\begin{tabular}{@{}llllll@{}}
\toprule
Configuration & Label precision & Label recall & Paint recall & GT pixels kept & Identity \\
\midrule
v7: all three off & 0.610 & 0.736 & 0.813 & 0.693 & 0.432 \\
+ lateral scope 5.5 m & 0.610 & 0.736 & 0.812 & 0.693 & 0.621 \\
+ same-side merge & 0.610 & 0.736 & 0.813 & 0.693 & 0.459 \\
+ appearance gate & 0.894 & 0.559 & 0.812 & 0.693 & 0.409 \\
v8: all three on (adopted) & 0.894 & 0.559 & 0.812 & 0.693 & 0.665 \\
\bottomrule
\end{tabular}
\end{table*}
\subsection{Training-composition dose response}
4 shows the dose response as it now stands, and the earlier curve did not survive its own check. That curve (base one seed, 4$\times$ six, 6$\times$ one, 8.5$\times$ six: 0.707/0.711/\textbf{0.763}/0.701) was measured under an earlier configuration whose identity values do not reproduce under the frozen protocol --- the same 6$\times$ checkpoint reads 0.763 there and 0.667 here --- so it was not comparable with the acceptance numbers and has been replaced. The figure now reports a \textbf{paired-seed re-evaluation} on the frozen protocol (seeds 42--47 for every dose whose checkpoints still exist; Table VI): 0$\times$ 0.658 (0.631--0.690), 4$\times$ 0.661 (0.648--0.674), 6$\times$ 0.665 (0.645--0.676). The paired per-seed differences against 0$\times$ are +0.003 (4$\times$) and +0.007 (6$\times$), with mixed signs (seed 45 is negative for both), so \textbf{the three doses are indistinguishable at six seeds} and the claim that 6$\times$ is the peak of the response does not survive this test. The 8.5$\times$ arm, whose checkpoints were removed in a storage pass, has a \emph{recorded} six-seed mean of 0.701 --- higher than 6$\times$, the opposite of the earlier curve's ordering. The delivered composition remains 6$\times$ because it was selected before the one-shot confirmation and changing it would require a newly sealed set; what the paper claims about the composition is now bounded to "no dose effect was demonstrated".
\begin{figure*}[!t]
\centering
\includegraphics[width=\textwidth]{../../logs/paper_build/figures/fig3_dose_response.png}
\caption{Dose response on the frozen protocol with paired seeds (42--47): the three re-evaluated doses are indistinguishable, and the recorded 8.5$\times$ value sits higher (Table VI).}
\label{fig:dose}
\end{figure*}
Line-masked training metrics cannot establish line learning. A separately frozen mechanical check
of the unattended learner is documented in the supplement; it is not an additional arm in this
dose experiment. The supplement sweeps the lateral
scope instead: tightening it from 6.0 m to 3.0 m raises identity (dev 0.582 $\rightarrow$ 0.713; limited-class
0.581 $\rightarrow$ 0.758) and removes candidates (477 $\rightarrow$ 166 and 166 $\rightarrow$ 127), while role agreement is flat --- the
lateral scope is an identity lever, not a role lever. Two further single factors were rejected:
stochastic weight averaging \cite{izmailov2018swa} over the last five epochs rescued no seed and a Tversky
loss-weight arm multiplied the false-positive frame rate on the negative side by 4.5 (11.5\% $\rightarrow$ 51.9\% of
eligible frames, +40.4 percentage points; 329 $\rightarrow$ 1730 false-positive pixels; maximum component 320 $\rightarrow$ 808 px),
which is why the delivered arm keeps the balanced loss. Across every recorded pixel arm the negative-side readings cluster rather than spread (supplement), which is why the negative-example dose had to be varied instead of a loss-weight knob. Extending the dose response to five metrics (supplement) shows that identity and role move with the dose while coverage and the off-road candidate fraction barely do. The availability cost of the adopted definition is also visible at the pairing stage itself: 3--4 paired frames lost per 117-frame run.
\begin{table}[!t]
\centering
\caption{Balanced-seed dose comparison on the frozen protocol (seeds 42--47 paired across doses). The three re-evaluated doses are indistinguishable; the recorded 8.5$\times$ value is higher than 6$\times$, so the earlier "6$\times$ is the peak" reading does not survive.}
\begin{tabular}{@{}lllll@{}}
\toprule
Dose & Identity mean & Identity range & Paired diff vs 0$\times$ & Role \\
\midrule
0$\times$ (base) & 0.658 & 0.631--0.690 & --- & 0.874 \\
4$\times$ & 0.661 & 0.648--0.674 & +0.0030 & 0.864 \\
6$\times$ (delivered) & 0.665 & 0.645--0.676 & +0.0069 & 0.883 \\
8.5$\times$ (recorded; checkpoints removed) & 0.701 & 0.655--0.744 & --- & --- \\
\bottomrule
\end{tabular}
\end{table}
\subsection{One-shot held-out confirmation}
The sealed set was read once, by the delivered candidate. Overall label-scope recall is \textbf{0.460} and
precision \textbf{0.763} (IoU 0.403; road IoU 0.633); per group, three of seven groups sit at recall
0.17--0.20 --- but not with the same precision: two of them hold precision 0.864 and 0.922 while the third
(a27) has precision \textbf{0.169} and sends 61\% of its predicted line pixels off the road, so it is a group the
candidate largely fails rather than a clean high-precision group. One group carries no line ground truth
(recall undefined, recorded as such), and the remaining groups range 0.43--0.63. On the negative side
22 of 105 eligible frames carry false line pixels (20.9\%), but those pixels are 0.084\% of eligible
pixels, the largest connected component is 1164 px, and no false candidate appears inside the control
region. 5 shows the confirmation. Two properties of this result deserve emphasis. First, it does
\textbf{not} reproduce the acceptance numbers of Section V-A: the independent set is harder and its group
composition differs. Composition is counted two ways and both are recorded: by package, 104 of the 200
frames come from certified line-free packages and 96 from line-bearing ones; by per-frame eligibility,
95 frames carry line truth and 105 are eligible negatives, the difference being frames whose own eligibility
differs from their package's label. The delivered candidate's claim is bounded by the numbers reported here
rather than by the acceptance table.
Second, the confirmation records that the instance-level probes (identity, role, reference coverage) were \textbf{not} run on
this set; they are UNKNOWN and are not reported as confirmed.
\begin{figure*}[!t]
\centering
\includegraphics[width=\textwidth]{../../logs/paper_build/figures/fig4_final_confirm.png}
\caption{One-shot final confirmation on the sealed road-disjoint set: per-group metrics (a) and negative-side diagnostics (b).}
\label{fig:confirm}
\end{figure*}
the supplement states the composition the confirmation ran on (200 frames across seven groups; digest and protocol hash printed in the title so the figure cannot be reattributed to another seal). the supplement is the audit that admitted it: every source package is compared against the 407 frame hashes already used, and the packages that showed overlap were dropped before sealing. Third, and relevant to the gap with Section V-A, the sealed
set's truth is \textbf{engine-certified} while the acceptance sets are \textbf{human-revised}: the two numbers are measured against
different truth rungs, not only different data.
\subsection{Closed loop}
\textbf{Definition change in closed loop.} With the baseline definition the closed-loop runs show centre
crossings of 8, 8, 8 and 13 frames; with the adopted definition no centre crossing was observed in any of the
four runs. That is a statement about four runs, not a rate: the same runs travel very little (stall fraction
0.93--1.00, see below), and zero events at near-zero exposure cannot be extrapolated. The adopted definition
pays for it with availability: the sensor lane-source rate falls from a median of 0.851 (range 0.231--0.975)
to 0.292 (0.097--0.524), a 66\% relative drop, and the paired-lane rate from 0.890 (0.809--0.986) to 0.528
(0.130--0.995). 6 shows the trade-off. Both effects are consistent with
the offline finding that the baseline keeps far and non-paint candidates whose pairings pull the car
across the lane centre.
\begin{figure*}[!t]
\centering
\includegraphics[width=\textwidth]{../../logs/paper_build/figures/fig8_closed_loop_tradeoff.png}
\caption{Closed-loop trade-off: centre crossings versus sensor lane-source rate, four runs per arm.}
\label{fig:closedloop}
\end{figure*}
\textbf{Acceptance failure and its mechanism.} The driving acceptance gate (goal + perception-led lateral
mode + strict arbitration, four runs) fails \textbf{0/4}: stall fraction 0.93--1.00, travelled 0.9--14.5 m
in 120 s, 76.5--88.8 m remaining to goal; the safety-side checks that were measured held
(0 collisions, 0 off-road frames, 0 reversing). the supplement shows the per-run hard-target matrix over the
recent runs. We then localised the blocker by measurement rather than conjecture. Over 12 town runs,
the frames whose reason is "planned vehicle body crosses lane boundary" number 512; in those frames
the car is stationary in 96\% of cases, the current body is inside the lane in 512/512, and the
planned crossing sits at a median of 2.50 m (509/512 within 2--4 m) --- i.e. just inside the 4 m
hard-stop threshold. The car is parked 0.86 m off the lane centre with its heading \textasciitilde{}11° off the lane
bearing, so the inflated body sweep of the path the planner actually generated crosses the boundary within
\textasciitilde{}2.5 m. What the record supports is that the \emph{generated} path is refused, repeatedly, in every frame of that
set; it does not establish that no dynamically feasible forward path exists, because the planner's candidate
set is not the set of all feasible paths. 7 shows the anatomy.
The qualified pose-sweep entry (source/config/exit/cleanup and exposure checked per pose) attempted
all nine poses: \textbf{one of nine was completed and eight were classified `unplaceable`} (`placement_rc=3`,
no equal-exposure drive measurement). The completed nominal baseline had 46 frames, 15 accepted-path
frames, a perception lane in 26, and a net displacement of 0.56 m; production placement shifted the
requested pose by 1.05 m, so it is a `production_alignment` observation rather than a fixed-actual-pose
experiment. The six non-baseline cases that were still outstanding were then rerun with the port-wait,
port-rotation and backoff entry; all six were again qualified as `unplaceable`. These records establish
a placement-availability failure, not a causal fixed-pose grid result: the requested pose, the actual
post-placement pose and the exposure qualification are not interchangeable. The old raw files remain
historical evidence and are not counted as qualified measurements. A fixed-actual-pose experiment would
require a placement contract that rejects pose drift before exposure, plus a clean restart between
poses. Within the one completed observation, lack of a perception lane frequently precedes path generation;
sweep rejection is a separate downstream event and is the principal reason in 8 of the 46 baseline
frames.
Per-run traces, arm-level distributions and lane-placement distributions are in the supplement.
Zero observed collisions, off-pavement frames and reversing in the four formal acceptance runs do not
establish safety, particularly under near-stationary operation. The body-footprint centre-crossing
criterion failed in all four runs (12--24 frames per run), even when the reference-point crossing metric
was low. Availability failures and safety-related body-boundary failures must both be reported.
\textbf{Four factors run, one proposal refuted before implementation.} We pre-registered five driving-layer
factors against this blocker; four were implemented and run, and the fifth (pavement re-centring) was
refuted by prior measurement before any code was written (Table III). A bounded "hold" window that lets a verified trajectory be reused past its
offer age raised motion (reuse 7 $\rightarrow$ 42 frames; travelled 3.15 $\rightarrow$ 8.5 m) but worsened the worst-case
body-crossing count (29 $\rightarrow$ 30). Pulling the planned path onto the perceived lane centre engaged in 39
frames and improved the worst case substantially (centre crossings 3 $\rightarrow$ 0; body crossings 84 $\rightarrow$ 39)
but worsened one right-side item (0 $\rightarrow$ 2). Removing the obstacle term from the on-pavement gate --- see
below --- doubled travelled distance (6.85 $\rightarrow$ 14.1 m) and halved the distance to goal (83.5 $\rightarrow$ 76.6 m)
but worsened body crossings (23 $\rightarrow$ 66) while the car moved twice as far. Swapping in the delivered
perception arm produced no separation on path availability (0.179 $\rightarrow$ 0.200, overlapping ranges) and a
worse worst case (36 $\rightarrow$ 49). All four implemented factors follow the pre-registered maximum rule and are reported as negative; none was
adopted, and the fifth remains untested by construction.
The arbitration histogram is the evidence behind the stall fraction: over 2443 settled frames (frames at
least 8 s into a run, pooled over the twelve town runs the figure names), "no drivable path" is the largest
single reason at \textbf{1383 frames}, followed by the planned body-sweep crossing at \textbf{512}, which together
account for the bulk of the stationary time.
\begin{table*}[!t]
\centering
\caption{The pre-registered driving-layer factors against the deadlock, with mechanism engagement,
motion read-outs and the pre-registered verdict. Four were implemented and run; the fifth was refuted by
prior measurement before implementation.}
\begin{tabular}{@{}llllll@{}}
\toprule
Factor & Mechanism engaged? & Travelled (median) & Stall (median) & Worst-case body-cross frames (control $\rightarrow$ factor) & Verdict \\
\midrule
Bounded hold window & yes (reuse 7 $\rightarrow$ 42 frames) & 3.15 $\rightarrow$ 8.5 m & 0.986 $\rightarrow$ 0.958 & 29 $\rightarrow$ 30 & rejected \\
Path re-centring & yes (39 frames) & 5.1 $\rightarrow$ 5.75 m & 0.984 $\rightarrow$ 0.980 & 84 $\rightarrow$ 39 & rejected (one item worse) \\
On-pavement gate (obstacle term removed) & yes (availability 0.07--0.52 $\rightarrow$ 0.36--0.90) & 6.85 $\rightarrow$ 14.1 m & 0.966 $\rightarrow$ 0.950 & 23 $\rightarrow$ 66 & rejected \\
Delivered perception arm & no (no separation) & --- & 0.979 $\rightarrow$ 0.995 & 36 $\rightarrow$ 49 & rejected \\
Pavement re-centring (as proposed) & not implemented & --- & --- & --- & refuted by prior measurement \\
\bottomrule
\end{tabular}
\end{table*}
\textbf{Why the lane was being dropped.} The on-pavement gate revokes the perception lane when the lane
centre does not lie on observed drivable surface. Instrumenting the gate's evidence over two
diagnostic runs (415 settled frames) shows that of 202 revoked frames, \textbf{all 202} were revoked by the
obstacle term (median 6 obstacle cells of \textasciitilde{}12 observed samples), while the number of samples that
were observed but outside the drivable mask is zero in 192 of them. The obstacle layer is populated
from a neural bird's-eye-view head, so on this link the "on-pavement" gate was in effect an
obstacle-prediction veto on the lane centre. the supplement shows the evidence and the availability funnel:
in the baseline configuration only 116 of 415 settled frames (28\%) have a perception lane and only
79 (19\%) a planner path.
\subsection{Real-time behaviour and a metric defect}
\textbf{Control cadence.} Per-frame telemetry shows the perception-and-planning tick costs 377--432 ms, of
which 355 ms is sensor reading (camera ring 139--162 ms, range image 216--239 ms) and 11--13 ms is
planning; the neural heads cost 3--5 ms. The renderer itself runs at \textasciitilde{}50 fps (20 ms per frame), so the
simulation is not the bottleneck: the control loop is paced at 2 Hz of wall clock by design and is in any
case tick-bound at \textasciitilde{}2.5 Hz. Between ticks a sub-step path re-runs the controllers, up to 15 sub-step
iterations per tick, but it is gated on a path having been chosen: in the stall regime (no path) no sub-steps
run and commands are issued every \textasciitilde{}0.5 s, while in an arm where paths exist 12 sub-steps per frame yield
about 31 command emissions per second. These are command-emission rates on the wall clock, not a claim about
control bandwidth or about the age of the sensor data each command is computed from. Table IV summarises. The practical consequence is that the perceived "stutter" of the
vehicle is a control-cadence artefact of the path-availability blocker, not a rendering or compute
failure.
\begin{table*}[!t]
\centering
\caption{Control cadence per configuration: tick cost, sub-step count and command-emission rate
(wall clock).}
\begin{tabular}{@{}llllll@{}}
\toprule
Run & Tick (ms) & Sensor read (ms) & Renderer frame (ms) & Sub-steps per frame & Effective command rate \\
\midrule
Baseline acceptance (no path) & 432 & 401 & 20 & 0 & \textasciitilde{}2.2 Hz \\
Lane-gate factor arm & 377 & 355 & 20 & 1 (max 5) & \textasciitilde{}5 Hz \\
Path-re-centring arm & 403 & \textasciitilde{}370 & 20 & 12 & \textasciitilde{}31 Hz \\
\bottomrule
\end{tabular}
\end{table*}
\textbf{T18 sensor replay correctness.} After the T17 cadence diagnosis, we froze the source snapshot and ran the T18 replay contract. The first Tech attempt is retained but invalid for equal-exposure comparison: its telemetry span is 18.208 s against the 24 s qualification minimum. A protocol-identical retry qualified. On the first valid warmup capture, the old and current range-processing paths produced \textbf{bitwise-equal published geometry and ray-hit outputs in 48 same-input calls}; the best repeated p95 fell from 172.647 ms to 132.420 ms. The scope is deliberately narrow: fresh-tracker single-frame CPU processing, not whole-runtime equivalence, multi-frame tracking, safety truth, actuator execution or a model promotion. Sensor capture time, independent pavement/paint truth, body coverage and actuator acknowledgement remain UNKNOWN; these are not treated as passes. See `T18_SENSOR_REPLAY_CORRECTNESS_RESULT_20261011.md`.
\textbf{A defect in the safety reading rule.} The safety item used to judge every driving factor is a frame
\emph{count} of body-centre crossings, and a factor arm can travel twice as far as its control. the supplement
quantifies the consequence: the same arm can halve its per-metre crossing rate while showing a larger
raw count. This is why, after the first two verdicts, we re-registered a per-distance reading
(crossings per 100 m) alongside the raw maximum; the earlier verdicts stand as given, and we flag the
rule defect as a methodological result rather than editing them.
\section{Discussion}
\textbf{What generalises.} Five artefacts from this programme seem reusable beyond the specific stack.
(i) The \textbf{provenance ladder with enforced credentials}: naming the truth rung of every label, and refusing
to let a command-line string promote one, made "which truth is this number measured against" an answerable
question. (ii) The \textbf{counting contract}: reporting every ratio with its denominator, and keeping frame-level
and instance-level counters separate, removed an entire class of "good number, empty pool" errors and made
per-scene failures visible instead of averaged away. (iii) **Protocol versioning with a hash bound to the
sealed set**: the mechanism is cheap and it caught a real mislabelling defect; without it, an evaluation
produced under one definition would have been attributed to another. (iv) **Appearance certification of
generated labels**: requiring that a claimed line be recoverable from the image turned "annotation noise"
from a tuning problem into an admission decision. (v) **Pre-registered single-factor A/Bs with a fixed
reading rule**, including an explicit refusal to treat unmeasured quantities as passes; this is what allowed
four negative results to be reported as results rather than as failures.
\textbf{What did not work.} The scope gap is not closable by engineering alone, and it is worth being precise
about why. About 30\% of annotated line pixels fail the appearance predicate the adopted definition uses;
the predicate reads RGB brightness and chroma only and computes no spatial contrast against the surrounding
pavement, so what the measurement establishes is that those pixels do not satisfy \emph{this} predicate. Whether
they are visually identifiable to a human, recoverable from context or from another modality, or learnable at
all are open questions that this work does not answer --- the strong form of the claim ("invisible, therefore
unlearnable") is not supported by the measurement. What is supported is the arithmetic: the paint scope
retains 69\% of annotated line pixels on the development pool, and the retained pixels are recalled at a
higher rate than the removed ones (0.813 versus 0.736 at the same checkpoints), so conditioning the recall on
the predicate is not a free improvement and is always reported alongside the full-label reading. Training on
the removed pixels remains untested rather than refuted.
The offline gates do not transfer to closed-loop acceptance automatically either: the same perception that
clears every offline gate leaves the car unable to move, for reasons (a static body-sweep deadlock) that no
offline metric expresses. Finally, availability and safety trade against each other along the definition
axis in a way that a single scalar cannot capture: in the four runs we have, the adopted definition was
accompanied by no observed centre crossing and by a 66\% relative drop in lane-source availability.
\textbf{Implications for evaluation practice.} Three practices would have changed our results materially.
Report the \emph{definition} alongside the metric, and make the definition a versioned artefact. Report
per-distance or per-opportunity normalisations for safety counts, because motion confounds raw
counts. And run one independent, one-shot confirmation whose composition differs from the
development pool; our confirmation's negative result is the single most informative measurement in
the programme.
\begin{table}[!t]
\centering
\caption{Configuration of the evaluated arms, read from the delivered candidate's checkpoint metadata. Every arm shares this configuration; the composition arms differ only in training-pool composition and the single-factor arms only in the named switch.}
\begin{tabular}{@{}ll@{}}
\toprule
Item & Value \\
\midrule
Architecture & `SegUNet`, a lightweight U-Net \cite{ronneberger2015unet}: three encoder blocks (3x3 convolutions, batch norm, 2x max-pooling), transposed-convolution decoder with skip connections, 1x1 head; channel width factor 1.0 \\
Parameters & 834,931 (counted from the delivered checkpoint) \\
Input & 536 x 403 RGB frame \\
Classes & background, asphalt, line \\
Class weights & 0.147 / 1.0 / 40.0 (line class up-weighted) \\
Loss & weighted cross-entropy + soft-clDice on the line channel (weight 1.0) + Tversky term \cite{salehi2017tversky} (alpha 0.3, beta 0.7, weight 1.0); the rejected arm of Section V-C differs only in the Tversky weights \\
Optimiser & Adam (beta 0.9/0.999, no weight decay), lr 1e-3 \\
LR schedule & CosineAnnealingLR, T_max 480 steps, stepped per iteration \\
Batch size & 4 \\
Training budget & 480 steps (40 per epoch; 12 epochs recorded across resumes) \\
Sampling & quota sampler over the full pool; per-run validation fraction 0.2 (158 train / 29 validation frames in the delivered arm) \\
Precision & AMP enabled; deterministic algorithms requested (CUDA cross-entropy has no deterministic implementation, which is recorded rather than hidden) \\
Initialisation & every arm starts from the same checkpoint (`t14_e1_promotable_20260927/baseline/seed42`, sha16 `b3a4dd5a82dd4096`); that checkpoint's own initialisation lineage is \textbf{not} fully recorded (`provenance_complete = False`), so the ancestry above it cannot be reconstructed from the stored metadata \\
Hardware & NVIDIA RTX 5070; torch 2.12.0.dev+cu128, CUDA 12.8, Python 3.10.11 \\
Matching (frozen) & a candidate matches engine truth when a truth line lies within 0.8 m lateral tolerance; reference availability requires annotation pixels on the candidate's own side; role agreement compares the candidate's left/right role with the engine line's role \\
Averaging & integer counters accumulated per frame, ratios computed once (never by averaging per-directory ratios); empty denominators are UNKNOWN with a stated reason, never 0 \\
Pools & human-revised anchor 5 directories / 75 frames; development pool 9 deduplicated scenes / 117 frames; limited-class pool 24 frames; sealed road-disjoint held-out set 200 frames \\
\bottomrule
\end{tabular}
\end{table}
\section{Limitations and Threats to Validity}
\textbf{No instance-level recall is claimed.} The counting contract reports coverage, identity and role as
conditional ratios over evaluated candidates; a missed marking on a frame where the perception produced no
candidate enters pixel-level recall only. Readers who need instance recall should read the contract's
definition (Section III-C) rather than infer one from $|R|/|C|$.
\textbf{A small human anchor.} The human-revised anchor is small (75 training frames) and the human-revised evaluation packs
cover 9 packages (117 frames), so scene diversity is limited; the sealed one-shot set is engine-certified rather than
human-revised, so its numbers are not directly comparable to the human-revised acceptance numbers. A larger human-revised
held-out set would strengthen the definitional argument and is left to future work. The programme deliberately added **no
new manual annotation** during the reported rounds, which bounds both the cost and the coverage.
\textbf{Single simulator, single map, and an unqualified second-map attempt.} The reported perception
results come from one simulator and one map. A controlled `west_coast_usa` batch has zero eligible
sites: the generator reports three isolated line sites and three not-applicable no-line sites.
Raw checks record zero line pixels at all sites. ROAD qualifies in 0/6 frames at five sites and
4/6 at the remaining site; every site's ROAD status is UNKNOWN. Alignment and visibility support
are also incomplete. A separate native collection there has 803 line-class
pixels over 30 forward-camera frames (about 27 per frame); the recorded `utah` forward view has
zero over 77 frames. These observations establish failure of this batch, not an engine-only
cause or the absence of usable truth across twenty maps. Camera pose, frame acquisition, material
visibility and class mapping remain competing explanations. The next automatic probe must
separate them before new data are admitted. The definitional effect's direction, but not its
magnitude, is reproduced on a second model (Section V-B). Driving results cover one town link.
\textbf{Six seeds, and unequal seed counts.} Acceptance uses six seeds; composition arms are reported with their
seed counts (base 1, 4$\times$ 6, 6$\times$ 1, 8.5$\times$ 6), and the 6$\times$ dose point rests on the single seed that was delivered.
The dose result is therefore not an effect size: the balanced-seed comparison has been run (Table VI)
and finds the three re-evaluable doses indistinguishable, while the recorded 8.5$\times$ value is higher than
6$\times$; the 8.5$\times$ checkpoints are gone, so that arm cannot be re-evaluated. Frame counts inside a pool are
highly correlated, so the seeds are the unit of variation we report and we do not treat pixels as independent
trials.
\textbf{The held-out set is consumed.} The one-shot confirmation was performed once, as designed, and the
set may not be read again; any future claim requires a newly sealed set. The instance-level probes
were not run on it and remain UNKNOWN.
\textbf{Artefact retention.} Per-epoch training snapshots and superseded run directories were removed in a
storage pass. Selected verdict files, manifests, telemetry traces, seal records and the delivered
checkpoint remain locally. Retained summaries do not recover removed training states or establish
complete reconstruction of every figure; the claim manifest records input scopes and remaining gaps.
\textbf{Driving results are from a research stack.} The closed-loop results describe our stack, not
automated driving in general; the acceptance gate is a project-internal target, and its thresholds
were fixed before the runs and never relaxed.
\section{Conclusion}
We set out to make lane perception for automated driving measurable on a small, explicitly credentialed truth base,
and to make its evaluation honest about definitions. A credentialed provenance ladder with four certification gates for its engine rung kept the machine-generated increment
trustworthy while the human anchor stayed small and unchanged; a counting contract gave every reported ratio an
explicit denominator; and a versioned, hash-bound protocol made the definition of "line" an artefact
rather than an assumption. The resulting measurements are informative precisely where they are
uncomfortable: the two natural definitions of a line each fail a different gate on the same data, and
the disagreement is a third of the annotated pixels; the adopted definition passes all offline gates
and records no reference-point centre crossings in four closed-loop runs, at half the availability,
while full-body crossing constraints still fail; and the one-shot
independent confirmation does not reproduce the development numbers, which we report as the headline
of that test. The closed-loop acceptance gate still fails, and we localised the blocker to a
measurable deadlock mechanism rather than a tuning shortfall. Every factor, including five negative
ones, was pre-registered and judged by rules fixed in advance; we consider that discipline, and the
recorded negative results it produced, the main methodological contribution of this work.
\section*{Reproducibility and Data Availability}
Reported measurements are tied to local recorded artefacts: per-run scorecards, per-frame telemetry
traces, run manifests, gate verdict files and sealed-set records. Figure scripts read these inputs;
schematics are drawn separately from code definitions. Current fig21 and fig22 have explicit paths
and pinned input digests; complete freezing of other historical figure inputs remains unfinished.
Training pools, checkpoints and the sealed set are not distributed with the public code. The
availability statement below describes the evidence archive and the audit of removed artefacts.
\section*{Nomenclature}
\begin{table}[!t]
\centering
\begin{tabular}{@{}ll@{}}
\toprule
Symbol & Meaning \\
\midrule
$P$ & eligible frames whose truth confirms a paint line \\
$C$ & evaluated perception candidates in $P$ frames \\
$C_{out}$ & candidates that fall on frames outside $P$ \\
$R$ & candidates in $C$ with an available reference on their own side \\
$M$ & matched instances, $M \subseteq R$ \\
$L$ & matched candidates with decidable predicted and reference roles, $L \subseteq M$ \\
$A$ & instances whose left/right role agrees, $A \subseteq L$ \\
coverage & $R/C$ (candidate-reference availability) \\
identity & $M/R$ (candidate identity) \\
role & $A/L$ (role agreement) \\
precision, recall & pixel level, reported in label scope and in paint scope \\
\bottomrule
\end{tabular}
\end{table}
\begin{table}[!t]
\centering
\begin{tabular}{@{}ll@{}}
\toprule
Abbreviation & Meaning \\
\midrule
FSD & full self-driving (perception-led driving stack) \\
BEV & bird's-eye view \\
IoU & intersection over union \\
TTC & time to collision \\
SWA & stochastic weight averaging \\
PLC & painted-line centring corrector \\
A/B & alternating control/factor comparison \\
UNKNOWN & measured quantity unavailable; never read as a pass \\
\bottomrule
\end{tabular}
\end{table}
\section*{Open problems and unfinished work}
The programme is deliberately reported with its unfinished parts in the same detail as its
results; this section lists them with the evidence that localises each one and the concrete
next experiment.
\textbf{1. The closed-loop acceptance gate is unmet.} Four runs of the acceptance configuration fail
0/4 on the benchmark's hard targets (stall fraction 0.93--1.00, travelled 0.9--14.5 m in 120 s,
76.5--88.8 m remaining to goal), while the measured safety items hold (0 collisions, 0 off-road
frames, 0 reversing). The blocker is localised to a static deadlock: a stationary car, off-centre by 0.86 m, whose inflated body
sweep crosses the lane boundary within a median 2.50 m, so the path the planner generates is refused and the
car cannot re-centre (Section V-E, Figs. 7 and the supplement). The record shows that the
generated path is refused, not that no feasible path exists; separating the two needs a controlled sweep over
frozen poses that stores the perception boundary, the candidate paths, the sweep geometry and the final
command for both executable and non-executable cases. The exploratory 3x3 probe has nine output files,
but is not qualified as a fixed-actual-pose, equal-exposure experiment (Section V-E). Its single-point
pose-set interpretation is withdrawn. \emph{Next experiment}: repair deadline and placement failure handling,
record actual post-placement poses, and repeat the invalid conditions under a predeclared pose contract
and exposure qualification. Only then judge the three factors that
engaged (bounded hold window, path re-centring, lane gate) as a single pre-registered \textbf{combined} arm
against the best single-factor arm, plus the plan-layer decisions that the programme may not take
unilaterally --- the planned-crossing hard threshold and the semantics of the on-pavement gate, both of
which are safety rules.
\begin{figure*}[!t]
\centering
\includegraphics[width=\textwidth]{../../logs/paper_build/figures/fig5_deadlock_anatomy.png}
\caption{Deadlock anatomy: the planned crossing sits just inside the 4 m threshold (a) and the car is stationary in 96\% of those frames (b).}
\label{fig:deadlock}
\end{figure*}
\textbf{2. Lateral-reference quality is the next perception-side lever.} The planner's path sits a
median 0.773 m from its own lane reference while the two layers' references agree exactly
(0.0 m gap), and the bounded centring corrector engages only where line evidence exists (39 of
805 settled frames). \emph{Next experiment}: raise line-evidence availability (a stronger line head,
or a fusion that does not depend on the sparse line class), or move the centring upstream so the
path is planned on the lane reference rather than corrected after the fact.
\textbf{3. Dose response: done for three of four doses, and it did not reproduce.} The paired-seed comparison
has been run (Section V-C, Table VI): 0$\times$/4$\times$/6$\times$ are indistinguishable at six seeds. What remains is the
8.5$\times$ arm, whose checkpoints were removed in the storage pass, and any dose above 8.5$\times$; both need training
runs under the frozen protocol rather than re-evaluation. If a dose effect is to be claimed at all, it
needs equal seed counts per dose on more than one map, with the seed as the unit of variation.
\textbf{4. Control cadence.} The perception-and-planning tick costs 377--432 ms, of which 355 ms is
sensor reading (camera ring plus range image) and 11--13 ms is planning; the control loop is paced
at 2 Hz by design and is tick-bound at \textasciitilde{}2.5 Hz. Sub-steps smooth the interval at up to 15 Hz but
are gated on a path having been chosen, so the stall regime is exactly where they do not run.
\emph{Next experiment}: cheaper sensor reads (ring resolution or round-robin mounts, range reuse under
the existing keep-alive) measured as a single factor with both the tick cost and the driving hard
targets as read-outs.
\textbf{5. Specification of safety metrics.} The driving factors were judged with a frame-count
maximum, and a factor arm can travel twice as far as its control; the supplement quantifies
how the same arm can halve its per-metre crossing rate while showing a larger raw count. A
per-distance reading has been pre-registered for new rounds; the earlier verdicts stand as
recorded, and the rule defect is reported as a methodological result rather than edited away.
\textbf{6. The observed label/appearance gap.} The current appearance predicate excludes about 30\% of
annotated line pixels. This describes the measured predicate and scope; it does not prove that
all local-contrast models must fail or that additional training inevitably causes hallucination.
A stronger line head or another sensor modality requires a new controlled comparison. Any change
to truth scope, including removing pixels from a recall denominator, must be pre-registered and
reported separately from a model improvement.
\textbf{7. Truth provenance and held-out capacity.} The existing human-revised anchor is small (75
training frames; 9 evaluation packages, 117 frames), and the engine-certified one-shot set is
now \textbf{consumed}. Further confirmation requires a newly sealed set. The next automatic collection
must pass independent channel, alignment and provenance checks, retain UNKNOWN for unsupported
truth, and freeze development and confirmation data separately. No new manual annotation is
planned; results from different truth rungs must remain separate.
\textbf{8. Second map: attempted, with the cause unresolved.} The controlled batch fails road, line
and alignment qualification; this does not isolate class mapping as the sole cause. The next
automatic probe should record a visible known material, camera pose, RGB/annotation agreement
and raw class counts, changing one factor at a time. Only independently qualified data may enter
learning or a new frozen evaluation. Until then, scope sensitivity remains a one-map finding,
with its direction reproduced on a second model.
\textbf{9. Instance vocabulary for harder geometry.} The counting contract's role vocabulary
(left/right, near/far) is defined for lane markings on a single carriageway; junctions, 3-D lane
estimation, and multi-lane role conventions are undefined, and the identity gate's failure mode
on the baseline definition is exactly a role/instance ambiguity rather than a pixel error.
\section*{Declarations}
\textbf{Funding.} This research received no external funding.
\textbf{Competing interests.} The author declares no competing interests.
\textbf{Author contributions.} The author is the sole author and performed every role:
conceptualisation, methodology, software, validation, formal analysis, investigation, data
curation, writing --- original draft, writing --- review and editing, visualisation, supervision.
\textbf{Ethics.} Not applicable: the work uses a commercial driving simulator and no human subjects.
The human annotation referenced in this paper was performed by the author on his own simulator
recordings.
\textbf{Use of AI tools.} AI coding assistants (DeepSeek and OpenCode) were used as development tools
during the implementation and documentation of this work; all experimental design, measurements,
verdicts and conclusions were produced and verified by the author, and every quantitative claim in
this paper is regenerated from recorded artefacts by the scripts listed below.
\textbf{Data and code availability.} What is public is the \textbf{code and the manuscript sources}:
`https://github.com/Qiongkura/BeamNG-autopilot` (branch `fix/round3-hardening-20260921`) contains
the evaluation stack, the protocol and counting-contract definitions, the figure generators, the
manuscript Markdown, `build_paper.py` and `references.bib`. What is \textbf{not} public is the recorded
evidence: the per-run scorecards, per-frame telemetry traces, run manifests, gate verdict files,
sealed-set records and checkpoints live under `logs/` in the author's working copy, which is not
tracked by that repository (the ignore rules exclude it). Every number in this paper is regenerated
from those local artefacts by the scripts named in Appendix A, and `CLAIMS_20261007.md` lists, for
each claim, the exact input file, protocol, seed set and scope it rests on, so the mapping can be
checked file by file. A public archive of the artefacts is \textbf{not yet available}; requests should go to the author. What
does exist is the packaged evidence set itself: `scripts/m5_paper_evidence_archive.py` collects the
declared numerical evidence, corresponding source and frozen plot inputs into one archive together
with a machine-readable manifest of paths, sizes and SHA-256 digests, under `logs/paper_release/` in
the author's working copy. That directory is excluded from the public repository by the same rule as
the rest of `logs/`. The manifest records the actual file count, missing inputs and digest violations;
an incomplete archive fails validation. It provides a checklist for requesting this evidence, but
excludes model weights and raw training frames and does not by itself support full training
reproduction. Per-epoch training snapshots and superseded run directories removed in a storage pass
are listed in a machine-readable audit file (`logs/_cleanup_20261006.txt`), and that removal is why
the balanced-seed dose comparison of Section V-C covers three of four doses.
\section*{Acknowledgements}
The author thanks the BeamNG.tech team for the simulator and its scripting interface, and
gratefully acknowledges his parents and friends for their support. He also thanks the developers of
DeepSeek and OpenCode, whose coding assistants were used as development tools during this work.
\section*{Appendix A: Recorded artefacts behind each results section}
\begin{table}[!t]
\centering
\begin{tabular}{@{}ll@{}}
\toprule
Results section & Artefact(s) \\
\midrule
V-A acceptance & `r2_verdict_v7_20261005.json`, `r2_verdict_v8_ADOPTED_20261005.json` \\
V-B definitional gap & `r2_verdict_v8_ADOPTED_20261005.json` (per-seed label and paint recall) \\
V-B boundary map & `r2_verdict_v8_paint_{nodilate_mean,dil1b,dil2}_20261005.json` \\
V-B appearance gate & `dev_pixel_{base,appgate}_20261004.json`, `r3b_pixel_{lat3,appgate}_20261004.json` \\
V-B scans & `dev_scan_{lateral,parallel}_20261004.json`, `r3b_scan_{lateral,parallel}_20261004.json` \\
V-B identity scopes & `t16_dual_scope_10seed.json` \\
V-C dose response & `dev_full_arms_20261005.json`, `dev_full_dose_arms_s4{3..7}.json` \\
V-C rejected factors & `r2_verdict_swa_20261005.json`, `dev9_beta08_v8d1_20261005.json` \\
V-D confirmation & `final_set_v5_20261005/seal/{final_set_seal.json,confirmation_base6x-seed42.json,final_set_ledger.jsonl}` \\
V-D composition audit & `final_pool_audit_20261005.json`, `collect_isolation_20260927.json` \\
V-E driving & `logs/fsd_benchmark/{scorecard,manifest,town}_*.json` (all runs of every A/B) \\
V-F timing & `timing_retest_20260927.json`, per-frame telemetry `frame_ms`/`tick_ms` fields \\
IV pool screening & `t16_r3_density_dev.json`, `t16_r3_pool{,b}_density_scenes.json`, `e1_negative_pool_20260927.json` \\
IV cost ledger & `gpu_ledger_machine.json` \\
\bottomrule
\end{tabular}
\end{table}

\begin{thebibliography}{99}
\bibitem{neven2018lanenet} Neven, Davy, De Brabandere, Bert, Georgoulis, Stamatios, Proesmans, Marc, Van Gool, Luc, “Towards End-to-End Lane Detection: An Instance Segmentation Approach,” in IEEE Intelligent Vehicles Symposium (IV), 2018, doi: 10.1109/IVS.2018.8500547.
\bibitem{pan2018scnn} Pan, Xingang, Shi, Jianping, Luo, Ping, Wang, Xiaogang, Tang, Xiaoou, “Spatial As Deep: Spatial CNN for Traffic Scene Understanding,” in AAAI Conference on Artificial Intelligence, 2018, doi: 10.1609/aaai.v32i1.12301.
\bibitem{zheng2022clrnet} Zheng, Tu, Huang, Yifei, Liu, Yang, Tang, Wenjian, Yang, Zheng, Cai, Deng, He, Xiaofei, “CLRNet: Cross Layer Refinement Network for Lane Detection,” in IEEE/CVF Conference on Computer Vision and Pattern Recognition (CVPR), 2022, doi: 10.1109/CVPR52688.2022.00097.
\bibitem{chen2026clrnetlight} 陈淑慧, “CLRNet车道线检测算法的一种轻量化方法,” 软件工程, 29, no. 9, 2026. [In Chinese]
\bibitem{huang2026bspline} 黄炜, 等, “基于B样条曲线的端到端车道线检测方法,” 长安大学学报(自然科学版), 46, no. 3, 118--129, 2026, doi: 10.19721/j.cnki.1671-8879.2026.03.008. [In Chinese]
\bibitem{jia2026gnn} 贾子厚, 等, “基于GNN-Transformer模型的车道线检测方法,” 工程设计学报, 33, no. 3, 2026. [In Chinese]
\bibitem{li2026xca} 李睿, 敖银辉, 陈新盛, 李桂伸, “基于Transformer和XCA注意力的车道线检测方法,” 电子测量技术, 49, no. 7, 2026, doi: 10.19651/j.cnki.emt.2519485. [In Chinese]
\bibitem{zhu2024parallel} 朱威, 欧全林, 洪力栋, 何德峰, “基于图像序列的车道线并行检测网络,” 模式识别与人工智能, 34, no. 5, 434--445, 2021. [In Chinese]
\bibitem{hu2026multiscale} 胡鹏, 刘平, 胡遥, “基于多尺度特征聚合的车道线检测方法,” 汽车工程学报, 2026. [In Chinese]
\bibitem{shuai2026deeplab} 帅强, 等, “基于改进DeeplabV3+模型的车道线检测方法,” 太原科技大学学报, 47, no. 3, 208--214, 2026, doi: 10.3969/j.issn.1673-2057.2026.03.006. [In Chinese]
\bibitem{chen2018deeplab} Chen, Liang-Chieh, Zhu, Yukun, Papandreou, George, Schroff, Florian, Adam, Hartwig, “Encoder-Decoder with Atrous Separable Convolution for Semantic Image Segmentation,” in European Conference on Computer Vision (ECCV), 2018, doi: 10.1007/978-3-030-01234-2_49.
\bibitem{lu2026symmetry} 陆军, 张澍森, 王磊, “复杂场景下基于对称先验与关键点预测的车道线检测,” 智能系统学报, 2026, doi: 10.11992/tis.202602011. [In Chinese]
\bibitem{chen2025threed} 陈小海, 等, “深度卷积神经网络架构下智能车3D车道线检测研究,” 现代电子技术, 48, no. 24, 2025. [In Chinese]
\bibitem{tusimple2017} TuSimple, “TuSimple Lane Detection Challenge,” 2017. [Dataset]
\bibitem{cordts2016cityscapes} Cordts, Marius, Omran, Mohamed, Ramos, Sebastian, Rehfeld, Timo, Enzweiler, Markus, Benenson, Rodrigo, Franke, Uwe, Roth, Stefan, Schiele, Bernt, “The Cityscapes Dataset for Semantic Urban Scene Understanding,” in IEEE Conference on Computer Vision and Pattern Recognition (CVPR), 2016, doi: 10.1109/CVPR.2016.350.
\bibitem{borkar2012groundtruth} Borkar, Amol, Hayes, Monson, Smith, Mark T., “A Novel Lane Detection System With Efficient Ground Truth Generation,” IEEE Transactions on Intelligent Transportation Systems, 13, no. 1, 365--374, 2012, doi: 10.1109/TITS.2011.2173196.
\bibitem{mccall2006video} McCall, Joel C., Trivedi, Mohan M., “Video-Based Lane Estimation and Tracking for Driver Assistance: Survey, System, and Evaluation,” IEEE Transactions on Intelligent Transportation Systems, 7, no. 1, 20--37, 2006, doi: 10.1109/TITS.2006.869595.
\bibitem{izmailov2018swa} Izmailov, Pavel, Podoprikhin, Dmitrii, Garipov, Timur, Vetrov, Dmitry, Wilson, Andrew Gordon, “Averaging Weights Leads to Wider Optima and Better Generalization,” in Conference on Uncertainty in Artificial Intelligence (UAI), 2018. [arXiv:1803.05407]
\bibitem{ronneberger2015unet} Ronneberger, Olaf, Fischer, Philipp, Brox, Thomas, “U-Net: Convolutional Networks for Biomedical Image Segmentation,” in Medical Image Computing and Computer-Assisted Intervention (MICCAI), 2015, doi: 10.1007/978-3-319-24574-4_28.
\bibitem{salehi2017tversky} Salehi, Seyed Sadegh Mohseni, Erdogmus, Deniz, Gholipour, Ali, “Tversky Loss Function for Image Segmentation Using 3D Fully Convolutional Deep Networks,” in Machine Learning in Medical Imaging (MLMI), MICCAI Workshop, 2017, doi: 10.1007/978-3-319-67389-9_44.
\end{thebibliography}
\end{document}
